%%=============================================================================
%% Backendarchitectuur en beveiliging
%%=============================================================================

\chapter{Backendarchitectuur en beveiliging}%
\label{ch:backendarchitectuur}
\label{ch:architectuur-implementatie}

Dit hoofdstuk bespreekt de backend van Forcebit Ticketing. De nadruk ligt op de gelaagde structuur, de verdeling van verantwoordelijkheden en de manier waarop tickettoegang gecontroleerd wordt.

\section{Gelaagde structuur}

De backend bestaat uit vier projecten. Elke laag heeft een duidelijke hoofdtaak:

\begin{itemize}[itemsep=0.2\baselineskip, topsep=0.25\baselineskip]
  \item \textbf{\texttt{1\_Api}}: ontvangt HTTP-verzoeken, leest JWT-claims, mapt requestdata en registreert afhankelijkheden.
  \item \textbf{\texttt{2\_Services}}: voert de use cases uit, gebruikt DTO's en bevat interfaces voor externe onderdelen zoals repositories en e-mail.
  \item \textbf{\texttt{3\_Persistence}}: bevat Entity Framework Core, repositories, configuraties en migraties.
  \item \textbf{\texttt{4\_Domain}}: bevat entiteiten, enums en domeinregels. Deze laag hangt niet af van de andere projectlagen.
\end{itemize}

De requestflow loopt meestal van API naar Services en daarna naar repositories. De API verwijst vooral naar de Persistence-laag om afhankelijkheden te registreren. Persistence implementeert de repositoryinterfaces uit Services. Daardoor kan de service-laag met abstracties werken, terwijl de concrete databankcode in de Persistence-laag blijft.

Deze opdeling ondersteunt Single Responsibility. Een controller hoeft niet te weten hoe een ticket uit MySQL wordt opgehaald. Een repository hoeft niet te beslissen of een klant een ticket mag sluiten. De Domain-laag blijft vrij van infrastructuurcode.

\begin{figure}[H]
  \centering
  \includegraphics[width=0.92\linewidth]{Backend_architecture.png}
  \caption[Backendlagen]{Backendlagen en belangrijkste projectafhankelijkheden.}
  \label{fig:backend-lagen}
\end{figure}

\section{Controllers en DTO's}

Controllers blijven bewust dun. Ze lezen routeparameters, requestdata, form-data en JWT-claims, maar voeren zelf geen ticketlogica uit. Daarna vertalen ze de binnenkomende data naar DTO's en geven die door aan de service-laag.

DTO's vormen hier de grens tussen HTTP en applicatielogica. Bij uploads ontvangt ASP.NET Core bijvoorbeeld een bestand als \texttt{IFormFile}. De service-laag werkt niet met dat HTTP-type, maar met \texttt{FileUploadRequest}. Zo blijven use cases bruikbaar zonder rechtstreeks afhankelijk te zijn van controllerdetails.

\section{Services als use cases}

\texttt{TicketService} is de belangrijkste use-caseklasse. Deze service maakt tickets aan, laadt ticketdetails, voegt berichten toe, wijzigt statussen, bewaart bijlagen en beslist wanneer een e-mail nodig is. De service stuurt de workflow aan, maar laat HTTP-antwoorden over aan controllers en databankqueries aan repositories.

Dit maakt de service verantwoordelijk voor de volgorde van acties, niet voor elke technische uitvoering. Bestanden opslaan gebeurt via \texttt{IFileStorageService}, e-mail via \texttt{IEmailService} en databanktoegang via repositories. Daardoor kan een onderdeel later vervangen worden zonder de volledige ticketflow te herschrijven.

\clearpage
\section{Domeinregels}

Businessregels rond tickets staan in de Domain-laag. Daardoor kunnen services dezelfde regels gebruiken zonder ze te kopiëren. De statusregel hieronder toont dat \texttt{New} alleen door ticketcreatie mag ontstaan. Een gebruiker kan een bestaand ticket dus niet terug naar \texttt{New} zetten.

\begin{listing}[H]
\begin{minted}[fontsize=\tiny,baselinestretch=0.82]{csharp}
if (newStatus == TicketStatus.New)
{
    return false;
}

if (userRole == UserRole.Admin)
{
    return true;
}

if (userRole != UserRole.Client || ticket.ClientId != userId)
{
    return false;
}

if (newStatus == TicketStatus.Closed)
{
    return ticket.Status != TicketStatus.Closed;
}

return newStatus == TicketStatus.Open &&
       ticket.Status == TicketStatus.Closed;
\end{minted}
\caption[Statuscontrole in TicketRules]{Statuscontrole in \texttt{TicketRules}.}
\label{lst:ticket-rules-status}
\end{listing}

De regel maakt ook het verschil tussen administrator en klant duidelijk. Een administrator mag een ticket naar \texttt{Open} of \texttt{Closed} zetten. Een klant mag alleen het eigen ticket sluiten of een gesloten ticket opnieuw openen. Daardoor blijft de betekenis van \texttt{New}, \texttt{Open} en \texttt{Closed} overal gelijk.

\section{Authenticatie en autorisatie}

Na login bevat het JWT-token de gebruikersidentiteit en rol. Controllers schermen routes af met rollen, terwijl de service-laag tickettoegang opnieuw controleert via domeinregels. Dat is nodig omdat een klant alleen eigen tickets mag zien, terwijl een administrator alle tickets mag beheren.

De frontend toont alleen acties die bij de rol passen, maar dat is niet voldoende als beveiliging. De backend blijft de beslissende laag. Ook wanneer iemand handmatig een API-aanroep probeert, controleert de service-laag opnieuw of de gebruiker de gevraagde actie mag uitvoeren.
